\documentclass[12pt]{article}
\usepackage[utf8]{inputenc}
\usepackage{graphicx}
\usepackage{amsmath}
\usepackage{amssymb}
\usepackage{mathrsfs}
\usepackage{pgfornament}
\usepackage[many]{tcolorbox}
\usepackage[margin = 0.5in]{geometry}
\usepackage{collectbox}
\usepackage{mdframed}
\usepackage{xcolor}

\newcommand\textbox[1]{%
  \parbox{.333\textwidth}{#1}%
}

\linespread{1.5}

\makeatletter
\newcommand{\mybox}{%
    \collectbox{%
        \setlength{\fboxsep}{2pt}%
        \fbox{\BOXCONTENT}%
    }%
}
\makeatother

\usepackage{listings}
\usepackage{color}

\definecolor{dkgreen}{rgb}{0,0.6,0}
\definecolor{gray}{rgb}{0.5,0.5,0.5}
\definecolor{mauve}{rgb}{0.58,0,0.82}

\lstset{frame=tb,
  language=R,
  aboveskip=3mm,
  belowskip=3mm,
  showstringspaces=false,
  columns=flexible,
  basicstyle={\small\ttfamily},
  numbers=none,
  numberstyle=\tiny\color{gray},
  keywordstyle=\color{blue},
  commentstyle=\color{dkgreen},
  stringstyle=\color{mauve},
  breaklines=true,
  breakatwhitespace=true,
  tabsize=3
}
\title{STS 532 HW III}
\begin{document}

%%
%% Title 
%%
\begin{center}
\vspace{-0.6cm}
\hrulefill \\
\vspace{-0.7cm}
\hrulefill

\vspace{-0.5cm}
\begin{flushleft}
\noindent\textbox{\pgfornament[width = 9mm, color = black]{39}\hfill}\textbox{\hfil \textbf{\begin{large}Mathematical Statistics\end{large}}\hfil}\textbox{\hfill \pgfornament[width = 9mm, color = black]{40}}
\end{flushleft}
\vspace{-5mm}
\begin{large}
\textbf{Homework III} \\
\end{large}

\vspace{5mm}

\begin{flushleft}
\noindent\textbox{\pgfornament[width = 9mm, color = black, symmetry = h]{39}\hfill}\textbox{\hfil \textbf{\begin{large}Nolan R. H. Gagnon\end{large}}\hfill}\textbox{\hfill \pgfornament[width = 9mm, color = black, symmetry = h]{40}}
\end{flushleft}

\vspace{-0.65cm}
\hrulefill \\
\vspace{-0.7cm}
\hrulefill
\end{center}

\begin{center}
\pgfornament[width = 75mm,
             color = black]{89}
\end{center}

\vspace{5mm}

%%
%% 8.12
%%
\noindent\mybox{\colorbox{lightgray}{\textbf{C \& B 8.12}}}\hspace{3mm}\hrulefill 

\noindent For samples of size $n = 1, 4, 16, 64, 100$ from a normal population with mean $\mu$ and known variance $\sigma^2$, plot the power function of the following LRTs. Take $\alpha = 0.05$.
\begin{flushleft}
\textbf{(a)} $H_0: \mu \leq 0$ versus $H_1: \mu > 0$ \\
\textbf{(b)} $H_0: \mu = 0$ versus $H_1: \mu \neq 0$
\end{flushleft}

\noindent\mybox{\colorbox{lightgray}{\textbf{Solution}}} \hspace{3mm}\hrulefill

\noindent \textcolor{blue}{\textbf{(a)}}  Following Example 8.3.3 (from Casella and Berger page 384), we see that the power function for this hypothesis test is 
\begin{align}
\beta(\mu) &= P\left(Z > c + \frac{\mu_0 - \mu}{\sigma/\sqrt{n}}\right) \\
&= P\left(Z > c - \frac{\mu}{\sigma/\sqrt{n}}\right). 
\end{align}

\noindent Our task now is to find the value of $c$ that forces the maximum probability of committing a Type I Error to be $\alpha = 0.05$. To this end, let $M_0 = (-\infty, 0]$ be the parameter space for the null hypothesis. If we force $$\beta(\sup M_0) = \beta(0)= 0.05,$$ then, due to the fact that $\beta(\mu)$ is an increasing function, we will have $\beta(\mu) \leq 0.05$ for all $\mu\in M_0$. But then we have guaranteed that the maximum probability of committing a Type I Error is $0.05$. Now, since we require $\beta(0) = P(Z > c) = 1 - P(Z \leq c) = 0.05$, we must choose $c = 1.645$ (obtained from the Standard Normal Probability Table). Therefore, the power function for this test is $$\beta(\mu) = P\left(Z > 1.645 - \frac{\mu}{\sigma/\sqrt{n}}\right).$$

Plots of this function for $n = 1, 4, 16, 64, 100$ are provided on the next page. Note that for graphing purposes, we took $\sigma = 1$.  

\begin{figure}[h!]
  \centering
  \includegraphics[width = 0.7\textwidth]{parta.png}	  \caption{Power function for the test $H_0: \mu \leq 0$ versus $H_1: \mu > 0$}
\end{figure}

From the plot above, it is clear that, as $n$ increases, the probability of committing a Type II Error decreases uniformly. 


\noindent \textcolor{blue}{\textbf{(b)}} The power function for this test is 
\begin{align}
\beta(\mu) &= P\left(Z \geq c - \frac{\mu}{\sigma/\sqrt{n}} \text{ or } Z \leq -c - \frac{\mu}{\sigma/\sqrt{n}}\right) \\
&= 1 - P\left(-c-\frac{\mu}{\sigma/\sqrt{n}} \leq Z \leq c - \frac{\mu}{\sigma/\sqrt{n}}\right) \\
&= 1 - P\left(Z \leq c - \frac{\mu}{\sigma/\sqrt{n}}\right) + P\left(Z\leq-c - \frac{\mu}{\sigma/\sqrt{n}} \right).
\end{align}
\noindent We require that
\begin{align}
\beta(0) &= 1 - P(Z\leq c) + P(Z\leq -c) \\
&= P(Z > c) + P(Z\leq -c) \\
&= 2P(Z\leq -c) \\
&= 0.05,
\end{align}
in order to control the maximum Type I Error probability. But this means $c = 1.96$ (from the Standard Normal Probability Table). Therefore, the power function of this hypothesis test is $$\beta(\mu) = 1 - P\left(Z \leq 1.96 - \frac{\mu}{\sigma/\sqrt{n}}\right) + P\left(Z\leq-1.96 - \frac{\mu}{\sigma/\sqrt{n}} \right).$$

\pagebreak

Plots of this function for $n = 1, 4, 16, 64, 100$ are provided below.

\begin{figure}[h!]
  \centering
  \includegraphics[width = 0.7\textwidth]{partb.png}	  \caption{Power function for the test $H_0: \mu = 0$ versus $H_1: \mu \neq 0$}
\end{figure}

Again, we see that, as $n$ increases, the probability of committing a Type II Error decreases uniformly.

\hfill$\blacksquare$

\vfill
\begin{center}
\pgfornament[width = 75mm,
             color = black]{89}
\end{center}
\begin{center}
\vspace{-0.6cm}
\hrulefill \\
\vspace{-0.7cm}
\hrulefill
\end{center}

\pagebreak

%%
%% 8.13
%%
\noindent\mybox{\colorbox{lightgray}{\textbf{C \& B 8.13}}}\hspace{3mm}\hrulefill 

\noindent Let $X_1, X_2$ be iid uniform$(\theta, \theta + 1)$. For testing $H_0: \theta = 0$ versus $H_1: \theta > 0$, we have two competing tests:

\begin{center}
$\phi_1(X_1): \text{ Reject } H_0 \text{ if } X_1 > 0.95,$ \\
$\phi_2(X_1, X_2): \text{ Reject } H_0 \text{ if } X_1 + X_2 > C$
\end{center}

\begin{flushleft}
\textbf{(a)} Find the value of $C$ so that $\phi_2$ has the same size as $\phi_1$.\\
\textbf{(b)} Calculate the power function of each test. Draw a well-labeled graph of each power function. \\
\textbf{(c)} Prove or disprove: $\phi_2$ is a more powerful test than $\phi_1$.\\
\textbf{(d)} Show how to get a test that has the same size but is more powerful than $\phi_2$.
\end{flushleft}

\noindent\mybox{\colorbox{lightgray}{\textbf{Solution}}} \hspace{3mm}\hrulefill

\noindent \textcolor{blue}{\textbf{(a)}} Let $\beta_1(\theta) = P_{\theta}(X > 0.95)$ be the power function for $\phi_1$.  The size of $\phi_1$ is
\begin{align}
\beta_1(0) &= P_0(X_1 > 0.95) \\
&= \int_{0.95}^1 1 \hspace{2mm} dx_1 \\
&= 0.05.
\end{align}

\noindent Now let $\beta_2(\theta) = P_{\theta}(X_1 + X_2 > C) = 1 - P_\theta(X_1 + X_2 \leq C)$ be the power function for $\phi_2$. The size of $\phi_2$ is $\beta_2(0) = 1- P_0(X_1 + X_2 \leq C)$. Examine the figure below.

\begin{figure}[h!]
  \centering
  \includegraphics[width = 0.6\textwidth]{hw3_813.png}	  \caption{Regions of Integration for calculating $P_0(X_1+X_2 \leq C)$}
\end{figure}

\noindent This figure shows the regions of integration for calculating $P_0(X_1 + X_2 \leq C)$.  If we require $\beta_2(0) = 0.05$, it must be the case that $P_0(X_1 + X_2 \leq C) = 0.95$.  Now, since the joint pdf of $X_1$ and $X_2$ is $f(x_1, x_2) = 1$ for $0 \leq x_1, x_2 \leq 1$ ($0$, otherwise), it follows that calculating $P_0(X_1 + X_2 \leq C)$ is the same as calculating the area of \textbf{$\text{R}_1$} or \textbf{$\text{R}_2$} (depending on the value of $C$). It is clear from the figure, however, that the area of \textbf{$\text{R}_1$} \textbf{cannot} be $0.95$.  The largest it can be is $0.5$, i.e., half the area of the unit square drawn in the figure. Thus, our region of integration must be \textbf{$\text{R}_2$}, which means $C$ must be some number between 1 and 2. So we have
\begin{align}
\beta_2(0) &= P_0(X_1 + X_2 > C) \\
&= 1 - P_0(X_1 + X_2 \leq C) \\
&= 1 - \left(1 - \int\limits_{C-1}^1\int\limits_{C-x_2}^1 1 \hspace{1mm} dx_1 dx_2\right) \\
&= \int\limits_{C-1}^1 \left[(1-C)+x_2\right] \hspace{1mm} dx_2 \\
&= (1-C) + 0.5 + (1-C)^2 - 0.5(1-C)^2.
\end{align}
Setting this equal to $0.05$ and rearranging yields the quadratic equation $$C^2 -4C +3.9 =0,$$ which has solutions $C_1 \approx 1.68377$ and $C_2 \approx 2.31623.$ Clearly $C$ cannot be greater than $2$, else $P_0(X_1 + X_2 > C) = 0$.  Thus, we choose $C = C_1 \approx 1.68377.$ 

\vspace{3mm}

\noindent \textcolor{blue}{\textbf{(b)}} The power function of $\phi_1$ is $\beta_1(\theta) = P_{\theta}(X_1 > 0.95).$  If $\theta > 0.95$, then $X_1$ is necessarily greater than $0.95$, so $P_\theta(X_1 > 0.95) = 1$.  If $\theta < -0.05$, then it is impossible for $X_1$ to be greater than $0.95$, so $P_\theta(X_1 > 0.95) = 0$. Finally, if $-0.05 < \theta < 0.95$, then 
\begin{align}
P_\theta(X_1 > 0.95) &= \int\limits_{0.95}^{\theta + 1} 1 \hspace{1mm} dx_1 \\
& = \theta + 1 - 0.95 \\
&= \theta + 0.05.
\end{align}
In summary, 
$$\beta_1(\theta) = 
\begin{cases}
0 & \text{ if } \theta < -0.05 \\
\theta + 0.05 & \text{ if } -0.05 < \theta < 0.95 \\
1 & \text{ if } \theta > 0.95.
\end{cases}
$$

A plot of $\beta_1(\theta)$ is given on the next page.

\begin{figure}[h!]
  \centering
  \includegraphics[width = 0.6\textwidth]{powerfunc1.png}	  \caption{Plot of $\beta_1(\theta)$}
\end{figure}
\pagebreak
Now $\beta_2(\theta) = 1 - P_\theta(X_1 + X_2 \leq 1.68377)$ will also be a piecewise function. Note that the smallest $X_1 + X_2$ can be is $2\theta$, and the largest it can be is $2\theta + 2$. If $\theta > 1.68377 / 2 = 0.8419$, then $P_\theta(X_1 + X_2 \leq 1.68377) = 0$ and $\beta_2(\theta) = 1$. If $\theta \leq \frac{1.68377 - 2}{2} = -0.1581$, then $P_\theta(X_1 + X_2 \leq 1.68377) = 1$, and $\beta_2(\theta) = 0$. The figure below will help us understand the behavior of $\beta_2(\theta)$ for $-0.1581 \leq \theta \leq 0.8419$. 

\begin{figure}[h!]
  \centering
  \includegraphics[width = 0.7\textwidth]{theta_regs.png}	  \caption{Regions of integration for $-0.1581 \leq \theta \leq 0.8419$}
\end{figure}

As can be seen in the figure, when $-0.1581 \leq \theta \leq 0.3419$,
\begin{align}
\beta_2(\theta) &= P_\theta(X_1+X_2 > 1.68377) \\
&= \int\limits_{1.68377-(\theta+1)}^{\theta+1}\int\limits_{1.68377-x_2}^{\theta+1} 1 \hspace{1mm} dx_1 dx_2 \\
&= \int\limits_{0.68377-\theta}^{\theta+1} (\theta - 0.68377 + x_2) \hspace{1mm} dx_2 \\
&= \left[(\theta-0.68377)x_2 +\frac{x_2^2}{2}\right]^{\theta+1}_{0.68377-\theta} \\
&= \frac{[1.68377 - 2(\theta+1)]^2}{2} \\
&= \frac{(0.31623 + 2\theta)^2}{2}.
\end{align}
For $0.3419 \leq \theta \leq 0.8419$,
\begin{align}
\beta_2(\theta) &= P_\theta(X_1 + X_2 >1.68377) \\
&= 1 - P_\theta(X_1+X_2\leq 1.68377) \\
&= 1 - \int\limits_{\theta}^{1.68377-\theta}\int\limits_{\theta}^{1.68377-x_2} 1 \hspace{1mm} dx_1 dx_2 \\
&= 1 - \int\limits_{\theta}^{1.68377-\theta} (1.68377 -x_2 - \theta) \hspace{1mm} dx_2 \\
&= 1 - \left[(1.68377-\theta)^2 - \frac{(1.68377-\theta)^2}{2}\right] + \left[(1.68377-\theta)\theta -\frac{\theta^2}{2}\right] \\
&= \frac{-0.83508 + 6.73508\theta - 4\theta^2}{2}.
\end{align}
In summary, 
$$
\beta_2(\theta) = 
\begin{cases}
0 & \text{ if } \theta \leq -0.1581 \\
\frac{(0.31623+2\theta)^2}{2} & \text{ if } -0.1581 \leq \theta \leq 0.3419 \\
\frac{-0.83508+6.73508\theta -4\theta^2}{2} & \text{ if } 0.3419 \leq \theta \leq 0.8419 \\
1 & \text{ if } \theta \geq 0.8419.
\end{cases}
$$
A plot of $\beta_2(\theta)$ is given on the next page.
\begin{figure}[h!]
  \centering
  \includegraphics[width = 0.6\textwidth]{beta_2.png}	  \caption{Plot of $\beta_2(\theta)$}
\end{figure}
\pagebreak

\noindent \textcolor{blue}{\textbf{(c)}}  A plot comparing $\beta_1(\theta)$ and $\beta_2(\theta)$ is provided below.

\begin{figure}[h!]
  \centering
  \includegraphics[width = 0.6\textwidth]{comp_power_funcs.png}	  \caption{Comparison of $\beta_1(\theta)$ and $\beta_2(\theta)$}
\end{figure}

\noindent From the plot above, it is clear that $\phi_2$ is not uniformly more powerful than $\phi_1$.  

\hfill$\blacksquare$

\vfill
\begin{center}
\pgfornament[width = 75mm,
             color = black]{89}
\end{center}
\begin{center}
\vspace{-0.6cm}
\hrulefill \\
\vspace{-0.7cm}
\hrulefill
\end{center}

\pagebreak

%%
%% 8.14
%%
\noindent\mybox{\colorbox{lightgray}{\textbf{C \& B 8.14}}}\hspace{3mm}\hrulefill 

\noindent For a random sample $X_1, X_2, \ldots, X_n$ of Bernoulli($p$) variables, it is desired to test
$$
H_0: p = 0.49 \text{ versus } H_1: p = 0.51.
$$
Use the Central Limit Theorem to determine, approximately, the sample size needed so that the two probabilities of error are both about 0.01. Use a test function that rejects $H_0$ if $\sum\limits_{i=1}^n X_i$ is large.

\vspace{3mm}

\noindent\mybox{\colorbox{lightgray}{\textbf{Solution}}} \hspace{3mm}\hrulefill

\noindent By the Central Limit Theorem, the distribution of $\frac{\sqrt{n}(\bar{X}_n - p)}{\sqrt{p(1-p)}}$ is approximately standard normal for large $n$.  But we can write
\begin{align}
\frac{\sqrt{n}(\bar{X}_n - p)}{\sqrt{p(1-p)}} &= \frac{\sqrt{n}}{\sqrt{p(1-p)}}\left(\frac{1}{n}\sum\limits_{i=1}^n X_i - p\right) \\
&= \frac{\sqrt{n}}{\sqrt{p(1-p)}}\left(\frac{\sum\limits_{i=1}^n X_i - np}{n}\right) \\
&= \frac{\sum\limits_{i=1}^n X_i - np}{\sqrt{np(1-p)}}.
\end{align}
Now, the power function for this hypothesis test is 
\begin{align}
\beta(p) &= P_p\left(\sum\limits_{i=1}^n X_i > c\right) \\ 
&= P_p\left(\frac{\sum\limits_{i=1}^n X_i - np}{\sqrt{np(1-p)}} > \frac{c - np}{\sqrt{np(1-p)}}\right) \\
&\approx P\left(Z > \frac{c-np}{\sqrt{np(1-p)}} \right) \\
&= 1 - P\left(Z \leq \frac{c-np}{\sqrt{np(1-p)}} \right), 
\end{align}
where $Z$ has the standard normal distribution. The probability of committing a Type I Error is $$\beta(0.49) = 1 - P\left(Z \leq \frac{c-0.49n}{\sqrt{n(0.49)(1-0.49)}} \right)$$ and the probability of committing a Type II Error is $$1 - \beta(0.51) = P\left(Z \leq \frac{c-0.51n}{\sqrt{n(0.51)(1-0.51)}}\right).$$ We want both of these probabilities to be approximately $0.01$, which gives us the two equations
$$P\left(Z \leq \frac{c-0.49n}{\sqrt{n(0.49)(1-0.49)}} \right) = 0.99$$
$$P\left(Z \leq \frac{c-0.51n}{\sqrt{n(0.51)(1-0.51)}}\right) = 0.01.$$
Using the Standard Normal Probability Table, these equations become
$$\frac{c-0.49n}{\sqrt{n(0.49)(1-0.49)}} = 2.33$$
$$\frac{c-0.51n}{\sqrt{n(0.51)(1-0.51)}} = -2.33.$$
The solution to this system of equations is $(c, n) \approx (6783.5, 13567)$.  Thus, a sample of size $n = 13567$ is required to achieve the specified error probabilities.

\hfill$\blacksquare$

\vfill
\begin{center}
\pgfornament[width = 75mm,
             color = black]{89}
\end{center}
\begin{center}
\vspace{-0.6cm}
\hrulefill \\
\vspace{-0.7cm}
\hrulefill
\end{center}

\pagebreak

%%
%% 8.15
%%
\noindent\mybox{\colorbox{lightgray}{\textbf{C \& B 8.15}}}\hspace{3mm}\hrulefill 

\noindent Show that for a random sample $X_1, X_2, \ldots, X_n$ from a $n(0, \sigma^2)$ population, the most powerful test of $H_0: \sigma = \sigma_0$ versus $H_1: \sigma = \sigma_1$, where $\sigma_0 < \sigma_1$, is given by
$$
\phi\left(\sum X_i^2\right) = 
\begin{cases}
1 & \text{ if } \sum X_i^2 > c \\
0 & \text { if } \sum X_i^2 \leq c.
\end{cases}
$$
For a given value of $\alpha$, the size of the Type I Error, show how the value of $c$ is explicitly determined. 

\vspace{3mm}

\noindent\mybox{\colorbox{lightgray}{\textbf{Solution}}} \hspace{3mm}\hrulefill

\noindent Let $f$ be the pdf of the $X_i$.  Suppose $$f(\textbf{x}|\sigma_1) > kf(\textbf{x}|\sigma_0),$$ where $k\geq 0$.  Then
\begin{align}
\frac{f(\textbf{x}|\sigma_1)}{f(\textbf{x}|\sigma_0)} &= \frac{\left(\frac{1}{\sqrt{2\pi}\sigma_1}\right)^n \exp\left\lbrace\frac{-1}{2\sigma_1}\sum\limits_{i=1}^n X_i^2 \right\rbrace}{\left(\frac{1}{\sqrt{2\pi}\sigma_0}\right)^n \exp\left\lbrace\frac{-1}{2\sigma_0}\sum\limits_{i=1}^n X_i^2 \right\rbrace} \\
&= \left(\frac{\sigma_0}{\sigma_1}\right)^n \exp\left\lbrace\left( \frac{1}{2\sigma_0} - \frac{1}{2\sigma_1}\right)\sum\limits_{i=1}^n X_i^2 \right\rbrace \\
&> k.
\end{align}
Multiplying by $(\sigma_0/\sigma_1)^n$ on both sides of the inequality and taking the natural logarithm yields
\begin{align}
\left( \frac{1}{2\sigma_0} - \frac{1}{2\sigma_1}\right)\sum\limits_{i=1}^n X_i^2 > \ln\left[\left(\frac{\sigma_1}{\sigma_0}\right)^n k\right].
\end{align}
Now, since $\sigma_0 < \sigma_1$, the quantity $\left( \frac{1}{2\sigma_0} - \frac{1}{2\sigma_1}\right)$ is positive, and dividing by it on both sides of the above inequality does not change the direction of the inequality symbol. So we have $$\sum\limits_{i=1}^n X_i^2 > \frac{\ln\left[\left(\frac{\sigma_1}{\sigma_0}\right)^n k\right]}{\left( \frac{1}{2\sigma_0} - \frac{1}{2\sigma_1}\right)} = c.$$ Therefore, if $f(\textbf{x}|\sigma_1) > k f(\textbf{x}|\sigma_0)$, then $\textbf{x}$ is in the rejection region. By similar reasoning, we can show that $\textbf{x}$ is not in the rejection region (i.e., $\sum\limits_{i=1}^n X_i^2 \leq c$) if $f(\textbf{x}|\sigma_1) < k f(\textbf{x}|\sigma_0).$ Thus, by the Neyman-Pearson Lemma, $\phi$ is the most powerful test of $H_0: \sigma = \sigma_0$ versus $H_1: \sigma = \sigma_1$.  

Now, for a given value of $\alpha$, we can calculate $c$ using 
\begin{align}
P_{\sigma_0}\left(\sum\limits_{i=1}^n X_i^2 > c\right) &= P\left(\sigma_0^2\sum\limits_{i=1}^n \left(\frac{X_i}{\sigma_0}\right)^2 > c\right) \\
&= P\left(\sigma_0^2Y > c\right) \\
&= P\left(Y > c / \sigma_0^2\right) \\
&= \alpha,
\end{align}
where $Y \sim \chi^2(n)$. Thus $c/\sigma_0^2 = \chi^2_{n, \alpha}$, or $$c = \chi^2_{n, \alpha}\sigma_0^2.$$

\hfill$\blacksquare$

\vfill
\begin{center}
\pgfornament[width = 75mm,
             color = black]{89}
\end{center}
\begin{center}
\vspace{-0.6cm}
\hrulefill \\
\vspace{-0.7cm}
\hrulefill
\end{center}

\pagebreak

%%
%% 8.19
%%
\noindent\mybox{\colorbox{lightgray}{\textbf{C \& B 8.19}}}\hspace{3mm}\hrulefill 

\noindent The random variable $X$ has pdf $f(x) = e^{-x}, x > 0$. One observation is obtained on the random variable $Y = X^\theta$, and a test of $H_0 : \theta = 1$ versus $H_1: \theta = 2$ needs to be constructed. Find the UMP level $\alpha = 0.10$ test and compute the Type II Error probability.

\vspace{3mm}

\noindent\mybox{\colorbox{lightgray}{\textbf{Solution}}} \hspace{3mm}\hrulefill

\noindent The pdf of $X$ indicates that it has an exponential distribution with parameter $\beta = 1$. Therefore, $Y$ must have a Weibull distribution with parameters $\gamma = \frac{1}{\theta}, \beta = 1$. The pdf of $Y$ is then 
$$g(y|\theta) =
\begin{cases}
\frac{1}{\theta}y^{1/\theta - 1}\exp\left\lbrace-y^{1/\theta} \right\rbrace & \text{ if } 0 \leq y < \infty \\
0 & \text{ otherwise. }
\end{cases}
$$  
By the Neyman-Pearson Lemma, a test $\phi$ with rejection region $R$ satisfying $\textbf{y} \in R$ if $g(\textbf{y}|2) > kg(\textbf{y}|1)$ (with $k\geq 0$) and power function $\beta(\theta)$ satisfying $\beta(1) = 0.10$ is an UMP level $\alpha = 0.10$ test of the hypotheses given in the problem. Now, because only one observation is taken on $Y$, the inequality $g(\textbf{y}|2) > kg(\textbf{y}|1)$ becomes $$\frac{1}{2}y^{-1/2}\exp\left\lbrace y-y^{1/2}\right\rbrace > k.$$ 
As can be seen in the figure below, the function on the left-hand side of the above inequality achieves a global minimum at $y = 1$. 

\begin{figure}[h!]
  \centering
  \includegraphics[width = 0.55\textwidth]{lhsineq.png}	  \caption{Plot of $f(y) = \frac{1}{2}y^{-1/2}\exp\left\lbrace y - y^{1/2} \right\rbrace$}
\end{figure}

\noindent Thus the information expressed in the inequality above can be re-expressed with the two inequalities $0 < y \leq k_1 \leq 1$ and $1 \leq k_2 \leq y$. If these inequalities describe the rejection region for the UMP test, then the size of this test is 
\begin{align}
\beta(1) &= P_1(Y \leq k_1 \text{ or } Y \geq k_2) \\
&= P_1(Y \leq k_1) + P_1(Y \geq k_2).
\end{align}
If $\theta = 1$, then the distribution of $Y$ is exponential with mean equal to $1$. Thus,
\begin{align}
P_1(Y \leq k_1) + P_1(Y \geq k_2) &= 1 - e^{-k_1} + e^{-k_2}.
\end{align}
We must find $k_1$ and $k_2$ such that $\beta(1) = 1 - e^{-k_1}+e^{-k_2} = 0.10$. We can do this using MATLAB.  A plot of $\beta(1)$ for different values of $k_1$ and $k_2$ is given below.

\begin{figure}[h!]
  \centering
  \includegraphics[width = 0.55\textwidth]{3d-plot.png}	  \caption{Plot of $1 - e^{-k_1}+e^{-k_2}$}
\end{figure}
\noindent From the plot, we can see that if we choose $k_1 = 0.082332$ and $k_2 = 3.864865$, then $1 - e^{-k_1}+e^{-k_2} \approx 0.1$.

For the $k_1$ and $k_2$ chosen above, the probability of making a Type II Error is
\begin{align}
1 - \beta(2) &= P_2(0.082332 \leq Y \leq 3.864865) \\
&= \int\limits_{0.082332}^{3.864865}\frac{1}{2}y^{-1/2}\exp\left\lbrace -y^{1/2} \right\rbrace \hspace{1mm} dy \\
&= \left[-\exp\left\lbrace -y^{1/2}\right\rbrace\right]_{0.082332}^{3.864865} \\
&\approx 0.61053.
\end{align}

\hfill$\blacksquare$
\vfill
\begin{center}
\pgfornament[width = 75mm,
             color = black]{89}
\end{center}
\begin{center}
\vspace{-0.6cm}
\hrulefill \\
\vspace{-0.7cm}
\hrulefill
\end{center}

\pagebreak

\end{document}
